\customchapter{ABSTRACT}

\begin{center}
    \large \vspace{-1.5cm} \textbf{TÍTULO DE LA TESIS EN INGLÉS}
\end{center}

This research describes the design, implementation, and integration of a teleoperated robotic system with haptic feedback, intended for performing superficial sutures in a controlled environment.
The system comprises two UR5 robotic arms, two Geomagic Touch haptic devices, and mechanical adaptations for the use of generic laparoscopic forceps.
The developed master-slave architecture enables an operator to remotely control the robot's end-effector while perceiving interaction forces through haptic feedback.

Various control strategies for trajectory tracking—including optimization-based control and sliding mode control (SMC)—were proposed and implemented, with their performance evaluated in terms of precision and stability.
Additionally, custom adapters were designed and fabricated for both Robotiq grippers and the haptic devices, ensuring proper interaction with the surgical instruments.

The system was validated through experimental tests using medical training kits, demonstrating its ability to perform superficial suturing tasks with precision, stability, and repeatability.
The results highlight the system's potential as an assistive or training tool and lay the groundwork for future research aimed at its application in real clinical settings.


\noindent \textbf{Keywords:}\\
\noindent Haptic control; Superficial suturing; UR5 robot; Optimization-based control; Sliding mode control.


